Выпускная квалификационная работа посвящена разработке генератора данных для
аналитического бенчмарка СУБД TPC-H. Актуальность работы определяется тем, что
традиционные средства формирования набора TPC-H ориентированы преимущественно на
вывод строковых файлов и не учитывают требований современной инфраструктуры
аналитических систем, использующей колонoчное представление данных, удаленные
объектные хранилища и встроенные сценарии генерации данных в оперативной
памяти.

Целью работы является разработка генератора Schengen, обеспечивающего
детерминированное формирование таблиц TPC-H с поддержкой параллельной
генерации, прямой записи в современные форматы хранения и отделением модуля
генерации от слоя сериализации.

Результатом работы является программная система на языке C++, выполняющая
генерацию таблиц TPC-H, формирование батчей Apache Arrow~\cite{apache-arrow},
запись в форматы Parquet~\cite{apache-parquet},
ORC~\cite{apache-orc}, Lance~\cite{lance-format} и
Vortex~\cite{vortex-format}, а также выгрузку результатов на локальные
и S3-совместимые хранилища. Отдельно реализована возможность
использования генератора как встраиваемого источника данных без промежуточной
записи на диск, что подтверждено интеграцией с СУБД
DuckDB~\cite{schengen-duckdb-repo}.
